\section{Institutional Background}
\label{sec:background}

\subsection{The July 2024 shutdown}
\label{sec:shutdown}

The shutdown followed escalating student protests over public-sector hiring quotas.
% TODO: ADD RELIABLE SOURCE FOR POLITICAL CONTEXT BEFORE FINAL SUBMISSION
The political context matters for identification because the decision to restrict connectivity was not randomly assigned. The empirical design therefore relies on the timing and severity of the network disruption, together with cross-country and within-Bangladesh comparisons, rather than assuming that the broader domestic environment was unchanged.

Available network measurements nevertheless provide a precise chronology. Cloudflare reports outages at major mobile operators beginning during the early hours of July 18, followed by a near-complete national loss of Internet traffic around 21:00 local time. Fixed broadband began to recover on July 23, initially as a limited restoration that prioritized sectors including banking, technology, exporters, outsourcing providers, and media. Some residential fixed networks returned later, while major mobile providers did not show substantial recovery until July 28 \citep{belson2024}. Internet Society Pulse independently records a national shutdown beginning July 18 and notes the gradual return of connectivity on July 23 \citep{isoc2024}.

The main treatment window in this paper is July 18--23 inclusive. This choice captures the six daily observations most clearly associated with the near-total national disruption while treating July 24 onward as the start of the restoration period in daily data. The distinction is necessarily approximate because restoration began during July 23 and differed across networks. For that reason, specifications based on alternative phase definitions and a regression discontinuity in time are used as robustness checks rather than relying on a single exact daily cutoff.

\begin{table}[htbp]
\centering
\small
\caption{Institutional timeline}
\label{tab:timeline}
\begin{threeparttable}
\begin{tabular}{@{}p{0.17\textwidth}p{0.43\textwidth}p{0.32\textwidth}@{}}
\toprule
Date & Connectivity status & Empirical treatment \\
\midrule
July 16--17 & Selective social-media/service blocking reported before the nationwide outage & Pre-treatment warning period; July 17 holiday handled in robustness \\
July 18 & Mobile networks disconnect; traffic falls toward zero nationally & Treatment begins \\
July 18--23 & Near-total nationwide Internet disruption & Main total-shutdown window \\
July 23--27 & Partial fixed-broadband restoration & Restoration / post-treatment transition \\
July 28 onward & Mobile connectivity begins to recover & Post-treatment \\
\bottomrule
\end{tabular}
\begin{tablenotes}[flushleft]
\footnotesize
\item \textit{Sources:} Cloudflare Radar reporting in \citet{belson2024} and Internet Society Pulse \citep{isoc2024}. Dates are operationalized at daily frequency for the empirical analysis.
\end{tablenotes}
\end{threeparttable}
\end{table}

\subsection{Why GitHub activity is informative about software production}
\label{sec:github}

Distributed version control separates code creation from code publication. Developers can create commits locally, but a push to a remote GitHub repository requires network connectivity. Other productive actions---opening or merging pull requests, reviewing code, closing issues, publishing releases, and commenting on collaborative work---are even more directly online. This distinction makes software development useful for studying an Internet shutdown: public GitHub activity measures an observable online margin, while commit metadata can reveal whether some code was created locally and uploaded later.

GH Archive records the public GitHub event stream and makes it available through hourly archives and Google BigQuery \citep{gharchive}. The data include events such as PushEvent, PullRequestEvent, PullRequestReviewEvent, IssuesEvent, IssueCommentEvent, and ReleaseEvent. For historical PushEvents, the event payload contains the number of commits and a list of commit objects, subject to a 20-commit timeline-event cap documented by GitHub. I account for this limitation in the mechanism analysis by defining a clean sample of fully observed pushes and by reporting robustness to progressively smaller push sizes \citep{githubevents}.
